% Davis (1963), "The Rotation of Eigenvectors by a Perturbation" --
% mathematical digest / proof skeleton.
%
% This file is not a transcription of the article.  It is a compact,
% independently written summary of the main mathematical mechanisms and
% proof ideas, with attribution to the original source:
%
%   Chandler Davis, "The Rotation of Eigenvectors by a Perturbation",
%   Journal of Mathematical Analysis and Applications 6, 159--173 (1963).
%   DOI: 10.1016/0022-247X(63)90001-5
%
% The default notation follows Davis's notation as closely as practical in
% LaTeX.  Switch \DavisOriginalNotationfalse to use a more modern notation.

\documentclass[11pt]{article}

\usepackage[T1]{fontenc}
\usepackage{lmodern}
\usepackage{microtype}
\usepackage{amsmath,amssymb,amsthm,mathtools}
\usepackage{mathrsfs}
\usepackage{xparse}
\usepackage[hidelinks]{hyperref}
\usepackage[margin=1.1in]{geometry}

\newif\ifDavisOriginalNotation
\DavisOriginalNotationtrue
% \DavisOriginalNotationfalse

\ifDavisOriginalNotation
  \newcommand{\Hsp}{\mathscr{H}}
  \newcommand{\Fsp}{\mathscr{F}}
  \newcommand{\Pinch}{\mathscr{C}}
  \newcommand{\Off}{\widetilde{\mathscr{C}}}
  \newcommand{\Comp}[1]{\widehat{#1}}
  \NewDocumentCommand{\blanknorm}{o}{\left\lVert\,\right\rVert\IfValueT{#1}{_{#1}}}
\else
  \newcommand{\Hsp}{\mathcal{H}}
  \newcommand{\Fsp}{\mathcal{F}}
  \newcommand{\Pinch}{\mathcal{D}}
  \newcommand{\Off}{\mathcal{D}^{\perp}}
  \newcommand{\Comp}[1]{#1^{\perp}}
  \NewDocumentCommand{\blanknorm}{o}{\left\lVert\cdot\right\rVert\IfValueT{#1}{_{#1}}}
\fi

\NewDocumentCommand{\norm}{om}{\left\lVert #2\right\rVert\IfValueT{#1}{_{#1}}}
\NewDocumentCommand{\inner}{mm}{\left(#1,#2\right)}
\newcommand{\tr}{\operatorname{tr}}
\newcommand{\spec}{\operatorname{spec}}
\newcommand{\Ran}{\operatorname{Ran}}
\newcommand{\Id}{I}
\newcommand{\eps}{\varepsilon}
\newcommand{\lam}{\lambda}
\newcommand{\hlam}{\lambda'}
\newcommand{\Pp}{P'}
\newcommand{\Phat}{\Comp{P}}
\newcommand{\Pphat}{\widehat{P'}}

\theoremstyle{definition}
\newtheorem{definition}{Definition}[section]
\newtheorem{observation}[definition]{Observation}
\newtheorem{result}[definition]{Result}
\newtheorem{argument}[definition]{Argument}

\title{Core Mathematical Arguments in Davis's Rotation-of-Eigenvectors Paper}
\author{Mathematical digest after Chandler Davis (1963)}
\date{}

\begin{document}
\maketitle

\begin{center}
\small
Original source: Chandler Davis, ``The Rotation of Eigenvectors by a Perturbation'',
\emph{Journal of Mathematical Analysis and Applications} 6, 159--173 (1963),
DOI: \href{https://doi.org/10.1016/0022-247X(63)90001-5}{10.1016/0022-247X(63)90001-5}.
This note is an independently written mathematical digest, not a transcription.
\end{center}

\section{Notation and the main problem}

Let \(A\) and \(H\) be hermitian operators on a finite-dimensional Hilbert
space \(\Hsp\).  The perturbation is \(A+H\).  Write the eigenvalues and
orthonormal eigenvectors of \(A\) as \(\lam_1,\ldots,\lam_n\) and
\(x_1,\ldots,x_n\), and write the corresponding data for \(A+H\) as
\(\hlam_1,\ldots,\hlam_n\) and \(x'_1,\ldots,x'_n\), when such a matching is
specified.

The two norms used most often are the operator norm \(\blanknorm\) and the
Frobenius norm \(\blanknorm[2]\).  Thus
\[
  \norm{B}=\sup_{\norm{x}=1}\norm{Bx},
  \qquad
  \norm[2]{B}^2=\tr B^2
\]
for hermitian \(B\).

Let \(P_1,\ldots,P_m\) be a complete orthogonal family of projectors, usually
spectral projectors of \(A\).  Davis's pinching operator is
\[
  \Pinch B = \sum_{j=1}^m P_j B P_j,
  \qquad
  \Off B = B-\Pinch B.
\]
The space \(\Fsp\) of hermitian operators, equipped with the Frobenius inner
product, makes \(\Pinch\) an orthogonal projector and \(\Off\) its complementary
projector.  The term \(\Off H\) is the part of the perturbation that mixes the
chosen spectral subspaces.

The problem is to bound how much the spectral subspaces of \(A\) rotate under
\(H\), preferably in terms of spectral gaps and norms of \(H\), especially
\(\norm{H}\), \(\norm[2]{H}\), and \(\norm[2]{\Off H}\).

\section{Canonical matching of subspaces}

Suppose \(P_1,\ldots,P_m\) and \(P'_1,\ldots,P'_m\) are two complete orthogonal
families of projectors with the same ranks, and assume no nonzero vector in
\(P_j\Hsp\) is annihilated by \(P'_j\).  There is a canonical unitary \(U\)
intertwining the two decompositions, obtained from the polar decomposition of
\(P'_jP_j\) on each summand:
\[
  U P_j
  = (P'_jP_jP'_j)^{-1/2} P'_jP_j
  = P'_j(P_jP'_jP_j)^{-1/2}P_j .
\]
It satisfies \(UP_j=P'_jU\).

\begin{observation}[Why this unitary is the right rotation]
Among unitaries \(W\) satisfying \(WP_j=P'_jW\) for every \(j\), the polar
unitary \(U\) minimizes the Frobenius displacement \(\norm[2]{\Id-W}\).  More
generally, it minimizes the unitarily invariant size of
\((\Id-W^*)(\Id-W))\) after pinching.  This justifies measuring the rotation of
spectral decompositions through quantities computed from \(U\).
\end{observation}

Choose an orthonormal basis in each \(P_j\Hsp\) consisting of eigenvectors of
\(P_jP'_jP_j\).  If \(x_i\) is one of these vectors, define its rotation angle by
\(\theta_i=\arccos\inner{Ux_i}{x_i}\).  Then
\[
  \norm[2]{\Off U}^2 = \sum_i \sin^2\theta_i,
  \qquad
  \norm[2]{\Id-U}^2 = \sum_i 2(1-\\cos\theta_i).
\]
Thus \(\norm[2]{\Off U}^2\) is the clean Frobenius ``sum of squared sines''
measure of total rotation.

\section{The single-cluster Rayleigh-quotient estimate}

Let \(P=E([-\beta,\beta])\) be a spectral projector of \(A\).  Assume that the
intervals \((-\beta-\gamma,-\beta)\) and \((\beta,\beta+\gamma)\) contain no
spectrum of \(A\).  Put \(\norm{H}=\delta<\gamma/2\), and let
\(P'=E'([-\beta-\delta,\beta+\delta])\) be the corresponding spectral projector
of \(A+H\).

\begin{result}[Single-cluster rotation bound]
For the complementary projector \(\Pphat=\Id-P'\),
\[
  \norm{\Pphat P} \le
  \frac{\beta+\delta}{\beta+\gamma-\delta}.
\]
\end{result}

\begin{argument}[Core proof]
Take \(x=Px\), \(\norm{x}=1\).  Since \(A\) is bounded by \(\beta\) on
\(P\Hsp\),
\[
  \inner{(A+H)^2x}{x}\le (\beta+\delta)^2.
\]
Decompose \(x=P'x+\Pphat x\).  On \(\Pphat\Hsp\), the operator \(A+H\) lies
outside the enlarged interval, so
\[
  \inner{(A+H)^2x}{x}\ge (\beta+\gamma-\delta)^2\norm{\Pphat x}^2.
\]
Combining the two inequalities gives the stated estimate.  The whole proof is
a Rayleigh-quotient comparison: the same quadratic form is bounded above by
what happens in the old spectral cluster and below by what happens outside the
new cluster.
\end{argument}

When \(\beta=0\), the cluster is a single eigenvalue, possibly with
multiplicity.  Retaining the discarded contribution on \(P'\Hsp\) gives the
sharper Frobenius estimate
\[
  \tr \Pphat P
  \le
  \frac{\tr H^2P-\tr(A+H)^2P'}{\gamma(\gamma-2\delta)}.
\]
Equivalently,
\[
  \norm[2]{\Pphat P}^2
  \le
  \frac{\norm[2]{HP}^2-\norm[2]{(A+H)P'}^2}{\gamma(\gamma-2\delta)}.
\]
The proof is the same comparison, summed over a basis of the old spectral
subspace chosen compatibly with the canonical unitary.

\section{Total rotation for simple spectrum}

Assume now that \(A\) has simple eigenvalues separated by at least \(\gamma\),
and \(\norm{H}=\delta<\gamma/2\).  Let \(U\) be the canonical unitary matching
one-dimensional spectral subspaces.  Summing the previous single-cluster
estimate gives a first global bound of the form
\[
  \norm[2]{\Id-U}^2
  \le
  \frac{2\norm[2]{H}^2}{(1+\cos\theta)\gamma(\gamma-2\delta)},
  \qquad
  \theta=\arcsin\frac{\delta}{\gamma-\delta}.
\]
This is useful, but the natural quantity is
\(\norm[2]{\Off U}^2=\sum_i\sin^2\theta_i\).

For each \(i\), define
\[
  \gamma_i = \min_{j\ne i}|\lam_i-\hlam_j|,
  \qquad
  (\gamma')^2=\min_i\{\gamma_i^2-(\lam_i-\hlam_i)^2\}.
\]
Then Davis's sharper global estimate is
\[
  (\gamma')^2\norm[2]{\Off U}^2
  \le
  \norm[2]{H}^2-
  \sum_i(\lam_i-\hlam_i)^2.
\]

\begin{argument}[Core proof]
For each old eigenvector \(x_i\), compare the same quantity in two ways:
\[
  \inner{(A+H-\hlam_i)^2x_i}{x_i}.
\]
On one side, split \(H=\Pinch H+\Off H\).  The cross term with \(\Off H\)
vanishes because \(\Pinch H\) is diagonal in the old eigenbasis, giving
\[
  \inner{(A+H-\hlam_i)^2x_i}{x_i}
  = (\lam_i-\hlam_i)^2 + \inner{(\Off H)^2x_i}{x_i}.
\]
On the other side, decompose \(x_i\) into its component in the new eigenspace
for \(\hlam_i\) and its component in the orthogonal complement.  The complement
is separated by at least \(\gamma_i\), so
\[
  \inner{(A+H-\hlam_i)^2x_i}{x_i}
  \ge
  (\lam_i-\hlam_i)^2 + (\gamma')^2\sin^2\theta_i.
\]
Summing over \(i\) yields the estimate.
\end{argument}

The important interpretation is that a fixed perturbation budget can be spent
on eigenvalue motion or on eigenspace rotation, but not independently on both.

\section{Lower bound on eigenvalue motion}

The sharper total-rotation theorem still contains the eigenvalue displacement
\(\sum_i(\lam_i-\hlam_i)^2\).  Davis removes it using a lower bound on how much
the diagonal part \(\Pinch H\) must move eigenvalues.

\begin{result}[Eigenvalue-change lower bound]
Suppose the eigenvalues of \(A+H\) are separated by at least \(\gamma>0\), and
\(\norm[2]{\Off H}\le \gamma/\sqrt2\).  Then
\[
  \sum_i(\hlam_i-\lam_i)^2
  \ge
  \norm[2]{\Pinch H}^2-\norm[2]{\Off H}^2.
\]
\end{result}

\begin{argument}[Geometric proof idea]
Work in the real Hilbert space \(\Fsp\) of hermitian matrices with Frobenius
norm.  Let \(B\) have the eigenvectors of \(A\) and the eigenvalues of
\(A+H\).  Let \(C=A+\Pinch H\).  The matrix \(C\) is the pinching of a matrix
unitarily equivalent to \(B\), so by the Schur-Horn convexity principle it lies
in the convex hull of the matrices obtained from \(B\) by permuting the new
eigenvalues.

The vertices of this convex polytope lie on a sphere.  The nearest nontrivial
permutation is obtained by swapping the pair of new eigenvalues with smallest
gap.  This gives a lower bound on the angle between \(-B\) and \(C-B\).  After
rewriting
\[
  (B-A)+(C-B)=\Pinch H,
  \qquad
  \Off(A+H)=\Off H,
\]
the angle inequality becomes the displayed lower bound.  Conceptually, the
diagonal part of the perturbation cannot be hidden: unless it is cancelled by
off-diagonal mixing, it must appear as eigenvalue displacement.
\end{argument}

Combining this result with the sharper total-rotation inequality gives the
clean off-diagonal control statement
\[
  (\gamma')^2\norm[2]{\Off U}^2 \le 2\norm[2]{\Off H}^2,
\]
under the corresponding smallness and gap hypotheses.  This is the central
structural message: eigenvector rotation is controlled primarily by the
part of \(H\) that is off-diagonal relative to the old spectral resolution.

\section{The sharp two-subspace estimate}

Finally specialize to two spectral subspaces separated by a gap.  After scaling,
assume \(A\) has no spectrum in \((-1,1)\), and let
\[
  P=E([1,\infty)), \qquad \Phat=E((-
\infty,-1]).
\]
Let \(x\) be an eigenvector of \(A+H\) with eigenvalue \(\hlam\ge 0\), and let
\(\theta\) be the acute angle between \(x\) and \(Px\).

\begin{result}[Sharp half-space rotation bounds]
If \(\norm{H}=\delta<1\), then
\[
  \sin 2\theta \le \delta.
\]
If in addition \(\Pinch H=0\), i.e. the diagonal blocks \(PHP\) and
\(\Phat H\Phat\) vanish, then the smallness assumption can be dropped and
\[
  \tan 2\theta \le \delta.
\]
Both constants are sharp in two-dimensional examples.
\end{result}

\begin{argument}[Two-dimensional compression]
Ignore the trivial cases \(Px=0\) and \(\Phat x=0\).  Compress to the span of
\(Px\) and \(\Phat x\), and use the orthonormal basis
\(Px/\norm{Px}\), \(\Phat x/\norm{\Phat x}\).  In this basis
\[
  x=\begin{pmatrix}\cos\theta\\ \sin\theta\end{pmatrix},
  \qquad
  A=\begin{pmatrix}a_1&0\\0&a_2\end{pmatrix},
  \qquad
  H=\begin{pmatrix}h_{11}&h_{12}\\\overline{h_{12}}&h_{22}\end{pmatrix},
\]
with \(a_1\ge 1\), \(a_2\le -1\).  The eigenvalue equation gives
\[
  \hlam=a_1+h_{11}+h_{12}\tan\theta
  =a_2+h_{22}+h_{12}\cot\theta.
\]
This forces the relevant off-diagonal entry to be real after the phase choice.
Rearranging yields
\[
  h_{12}(\cot\theta-\tan\theta)=a_1-a_2+h_{11}-h_{22}.
\]
Using \(a_1-a_2\ge2\) and the norm bound on the \(2\times2\) compression of
\(H\) gives the inequality \(\sin2\theta\le\delta\).  If
\(h_{11}=h_{22}=0\), the same displayed equation immediately gives
\(\tan2\theta\le\delta\).
\end{argument}

\section{Distilled logical dependencies}

\begin{enumerate}
  \item Match old and new spectral projectors with the polar-decomposition
  unitary \(U\).
  \item Measure rotation by \(\norm[2]{\Off U}^2=\sum_i\sin^2\theta_i\).
  \item For one cluster, compare the same Rayleigh quotient above and below to
  bound escape from the corresponding new cluster.
  \item Sum one-cluster bounds to control total rotation.
  \item Refine the total estimate by separating eigenvalue motion from
  eigenspace rotation.
  \item Prove that the diagonal part \(\Pinch H\) must contribute to eigenvalue
  motion, up to an error controlled by \(\Off H\).
  \item Conclude that the actual eigenspace rotation is controlled by the
  off-diagonal part \(\Off H\), not merely by the full perturbation norm.
  \item In the two-subspace case, compress to dimension two to get sharp
  trigonometric bounds.
\end{enumerate}

\end{document}
